\documentclass[12pt]{article}
\usepackage[T1]{fontenc}
\usepackage[utf8]{inputenc}
\usepackage{lmodern}
\usepackage{textcomp}
\usepackage{enumitem}
\usepackage{geometry}
\geometry{margin=1in}
\begin{document}
\begin{center}
Answer Key - Health Sciences - K-12 Level Assessment 6 \
\small (Answers are ordered exactly as in the question paper.)
\end{center}

\section*{Multiple Choice}
\begin{enumerate}[label=\arabic*.]
\item \textbf{Answer:} A
\\ \textit{Options:} A) N-3 fatty acids; B) Phytochemical; C) N-6 fatty acids; D) Long chain saturated fats
\\ \textit{Note:} N-3 fatty acids, particularly those found in fish, have been shown in various studies to support brain health and may help reduce the risk of cognitive decline associated with aging.
\item \textbf{Answer:} B
\\ \textit{Options:} A) Vitamin C; B) Dietary fibre; C) Alcohol; D) Oestrogen
\\ \textit{Note:} Dietary fiber is linked to a reduced risk of colorectal cancer because it promotes healthy digestion, enhances bowel regularity, and may help in the elimination of potential carcinogens from the digestive tract.
\item \textbf{Answer:} D
\\ \textit{Options:} A) 60 g glucose per hour; B) 60 g glucose plus fructose per hour; C) 90 g glucose per hour; D) 90 g glucose plus fructose per hour
\\ \textit{Note:} Elite athletes should ideally consume 90 g of glucose plus fructose per hour during prolonged high-intensity exercise because this combination enhances carbohydrate absorption and utilization, helping to maintain energy levels and improve performance over extended periods of exertion.
\item \textbf{Answer:} B
\\ \textit{Options:} A) apoA, apoB, apoC, apoD, and apoE; B) apoA, apo(a), apoB, apoC and apoE; C) apoA, apoB, apoC, apo E, and apoL; D) apoB, apoC, apoD, apoE and apoM
\\ \textit{Note:} The correct answer identifies the five main series of apoproteins-apoA, apo(a), apoB, apoC, and apoE-because these are the key proteins involved in lipid transport and metabolism, each playing distinct roles in the body's lipid profile and overall health.
\item \textbf{Answer:} B
\\ \textit{Options:} A) Liver glycogen; B) Muscle glycogen; C) Intramuscular lipid; D) Adipose tissue lipid
\\ \textit{Note:} Muscle glycogen is the primary source of energy during moderate to high intensity exercise because it is readily available for rapid breakdown into glucose, which is then used to produce ATP, the energy currency of the body.
\end{enumerate}

\section*{True / False}
\begin{enumerate}[label=\arabic*.]
\setcounter{enumi}{5}
\item \textbf{Answer:} False
\\ \textit{Options:} A) True; B) False
\\ \textit{Note:} Older adults are more likely to experience vitamin B$_{12}$ deficiency due to decreased absorption in the digestive system and dietary restrictions that limit their intake of B$_{12}$-rich foods.
\item \textbf{Answer:} True
\\ \textit{Options:} A) True; B) False
\\ \textit{Note:} Vitamin A is a crucial component of rhodopsin, a pigment in the retina that enables the eyes to detect light in low-light conditions, thus making it essential for vision in dim light.
\item \textbf{Answer:} True
\\ \textit{Options:} A) True; B) False
\\ \textit{Note:} Eosinophilic oesophagitis is characterized by an accumulation of eosinophils in the esophagus, leading to inflammation and swelling that can result in narrowing and obstruction of the esophagus.
\item \textbf{Answer:} True
\\ \textit{Options:} A) True; B) False
\\ \textit{Note:} High nutrient requirements for tissue turnover do not contribute to vitamin or mineral deficiencies in older people because it is the overall decline in dietary intake and absorption, rather than tissue turnover needs, that primarily leads to these deficiencies.
\item \textbf{Answer:} False
\\ \textit{Options:} A) True; B) False
\\ \textit{Note:} Vitamin A does not play a role in the activation of vitamin D receptors; rather, vitamin D itself is responsible for activating its own receptors in the body.
\end{enumerate}

\section*{Long Form Response}
\begin{enumerate}[label=\arabic*.]
\setcounter{enumi}{10}
\item \textbf{Answer:} A vegan diet can lead to potential nutritional deficiencies, particularly in Vitamin B$_{12}$, which is crucial for overall health. Vitamin B$_{12}$ is essential for the production of red blood cells, the maintenance of the nervous system, and the synthesis of DNA. Since B$_{12}$ is primarily found in animal products, vegans must be mindful to include fortified foods or supplements to meet their dietary needs. Without adequate B$_{12}$, individuals may experience fatigue, weakness, and neurological issues. Therefore, it is important for those following a vegan diet to ensure they are getting enough Vitamin B$_{12}$ to maintain their health and well-being.
\\ \textit{Note:} The answer correctly identifies Vitamin B$_{12}$ as a critical nutrient that is predominantly found in animal products, emphasizes its essential functions in the body, and highlights the need for vegans to seek alternatives to prevent deficiencies and associated health issues.
\item \textbf{Answer:} Maintaining specific blood pressure and lipid levels is crucial for adults with diabetes to prevent cardiovascular disease. For optimal heart health, blood pressure should be kept below 130/80 mmHg, triglycerides should be less than 150 mg/dL, and LDL cholesterol should remain under 100 mg/dL. High blood pressure and elevated lipid levels can lead to serious complications like heart attacks and strokes, which are more common in individuals with diabetes. By achieving these targets, individuals can significantly reduce their risk of cardiovascular problems and improve their overall health, leading to a better quality of life. Regular monitoring and lifestyle changes, such as a healthy diet and exercise, are essential in achieving these goals.
\\ \textit{Note:} correct because it emphasizes that maintaining specific blood pressure and lipid levels is essential for preventing cardiovascular complications in adults with diabetes, as these conditions can significantly increase the risk of heart disease and stroke.
\end{enumerate}

\vfill
\noindent\textit{End of Answer Key}
\end{document}